\section{Visuelles Redesign}\label{sec:visuelles-redesign}

Im Zuge des Redesigns wurde die Benutzeroberfläche nicht nur optisch modernisiert, sondern auch funktional gestrafft.
Ziel war es, eine ansprechendere, intuitivere und visuell klarere Benutzererfahrung zu schaffen.


\subsection{Ausgangszustand}\label{subsec:visuelles-redesign-ausgang}

Die ursprüngliche Version bestand aus separaten Seiten für Login und Registrierung sowie einer simplen, eher funktional ausgerichteten Struktur mit Sidebar-Navigation. Die Gestaltung war insgesamt sehr schlicht gehalten, mit wenigen farblichen Akzenten und begrenzter visueller Dynamik.

\subsection{Überarbeitung}\label{subsec:visuelles-redesign-overhaul}

Die neue Gestaltung setzt auf ein aufgeräumtes, lebendigeres Gesamtbild mit optimierter Benutzerführung und einem stärkeren Fokus auf Nutzerfreundlichkeit:

\begin{itemize}
    \item \textbf{Zusammenführung der Authentifizierung:} Anstelle separater Seiten für Login und Registrierung wurden beide Funktionen in die zentrale Startseite integriert. Per Klick kann nahtlos zwischen beiden Ansichten gewechselt werden, wobei der jeweilige Formularblock animiert verschoben wird. Dies schafft ein dynamisches und modernes Nutzungserlebnis.
    \begin{figure}[H]
        \centering
        \includegraphics[width=0.6\textwidth]{loginsignup_old}
        \caption{Ursprünglicher Login-Bereich (Praxismodul II)}
        \label{fig:login-old}
    \end{figure}
    \begin{figure}[H]
        \centering
        \includegraphics[width=0.6\textwidth]{loginsignup_new}
        \caption{Neuer Login-Bereich}
        \label{fig:login-new}
    \end{figure}

    \item \textbf{Top-Navigation statt Sidebar:} Die ehemals verwendete Sidebar wurde entfernt, da sie für die aktuelle Funktionalität nicht mehr notwendig war. Stattdessen wurde eine schlanke Navigationsleiste am oberen Seitenrand eingeführt, die eine klare Struktur bietet und gleichzeitig mehr Platz für den Hauptinhalt schafft.
    \begin{figure}[H]
        \centering
        \includegraphics[width=0.6\textwidth]{collectionPage_old}
        \caption{Sidebar (Praxismodul II)}
        \label{fig:sidebar-old}
    \end{figure}
    \begin{figure}[H]
        \centering
        
\includegraphics[width=0.6\textwidth]{sidebar_new}
        \caption{Topleiste als Sidebar-Ersatz}
        \label{fig:sidebar-new}
    \end{figure}

    \item \textbf{Gestalterische Aufwertung:} Die gesamte Oberfläche wurde farblich aufgefrischt und mit einem dezenten Kachelhintergrund versehen, um dem Design mehr Charakter und visuelle Tiefe zu verleihen. Die neue Farbgebung wirkt freundlicher und einladender.
    \begin{figure}[H]
        \centering
        \includegraphics[width=0.6\textwidth]{visuals_new_1}
        \caption{Beispiel Design Overhaul}
        \label{fig:visuals-new-1}
    \end{figure}
 \begin{figure}[H]
                 \centering
                 \includegraphics[width=0.6\textwidth]{visuals_new_2}
                 \caption{ Beispiel Landing Page Overhaul}
                 \label{fig:visuals-new-2}
    \end{figure}
    \item \textbf{Erweiterung der Funktionen:} Als Ergänzung wurde ein Logout-Button sichtbar integriert, was den Bedienkomfort erhöht und eine klarere Nutzersteuerung ermöglicht.
\end{itemize}

Durch die Neugestaltung wirkt die Oberfläche nicht nur frischer und moderner, sondern auch strukturierter.
Die Nutzerführung ist klarer, unnötige Elemente wurden entfernt und durch gezielte visuelle Anpassungen wirkt die Anwendung insgesamt stimmiger. Auch kleine Details wie Animationen und der neue Hintergrund sorgen dafür, dass die Seite lebendiger und einladender erscheint.
Die überarbeitete Navigation und der sichtbare Logout-Button runden die Nutzererfahrung sinnvoll ab.